% Drawing 31/05-2 -- folder 31, batch 05, archive page 95 (PDF page 96).
% Transcription: à gauche du lemme 3, un croquis : un cercle traversé de deux diamètres qui se coupent en un point marqué, le vertical épais ; « $V(y)$ » à droite, « $V(x)$ » au-dessous
% Kind: circle with two diameters V(x), V(y). Method 2 (schematic redraw from the facsimile). Status: ok.
% Notes: Ink, the circle doubled in pencil (drawn once). Heavy near-vertical diameter and thin near-horizontal one through the marked point; V(y) at the right of the horizontal one, V(x) under the vertical one. Redrawn 2026-11-12 (issue 128): to check again.
% Coordinates: 1 unit = 1 mm = 1/100 of the crop width (bbox in page fractions: [0.12, 0.62, 0.36, 0.77]).
% Redrawn by Claude (claude-fable-5-2), issue #128, 2026-11-12; not checked by a human; not compiled (no TeX in the session).
% Rights: reproduces part of the work -- behind RENDER_TRANSCRIPTS, never published by itself.
\documentclass[tikz,border=3pt]{standalone}
\usepackage{amsmath,amssymb}
\usetikzlibrary{patterns,arrows.meta}
\tikzset{pencil/.style={black!45}, struck/.style={black!35}, mpt/.style={circle,draw,inner sep=1.3pt}, pt/.style={circle,fill,inner sep=0.9pt}, opt/.style={circle,draw,fill=white,inner sep=1.1pt}}
\begin{document}
\begin{tikzpicture}[x=1mm,y=1mm,line cap=round,line join=round,font=\scriptsize]
\draw plot[smooth, tension=0.6] coordinates {(69.5,-34) (69.1,-38) (68.1,-41.8) (66.4,-45.4) (64,-48.7) (61.1,-51.5) (57.8,-53.7) (54,-55.4) (50.1,-56.5) (46,-56.8) (41.9,-56.5) (38,-55.4) (34.3,-53.7) (30.9,-51.5) (28,-48.7) (25.6,-45.4) (23.9,-41.8) (22.9,-38) (22.5,-34) (22.9,-30) (23.9,-26.2) (25.6,-22.6) (28,-19.3) (30.9,-16.5) (34.2,-14.3) (38,-12.6) (41.9,-11.5) (46,-11.2) (50.1,-11.5) (54,-12.6) (57.8,-14.3) (61.1,-16.5) (64,-19.3) (66.4,-22.6) (68.1,-26.2) (69.1,-30) (69.5,-34)};
\draw[very thick] (43.5,-11.5) -- (44.7,-33.3) -- (46,-57);
\draw (24,-35.5) -- (44.7,-33.3) -- (70,-32.5);
\node[pt] at (44.7,-33.3) {};
\node[inner sep=1pt] at (86,-32) {\small $V(y)$};
\node[inner sep=1pt] at (41,-66) {\small $V(x)$};
\end{tikzpicture}
\end{document}
